\section{Discussion}
\label{sec:discussion}

A reduced whole-cell model with a real ODE--stochastic boundary can be inferred from single-cell and
bulk data, and its posterior predictions carry part of the way. On the reduced \syn{} model the
mixed route gives one joint posterior over all \pn{kim92.ncoord} free parameters for about
\pn{kim92.cpuh} CPU-hours (\cref{sec:reduced}). The data reach every parameter, and the
posterior reproduces its own data. Pushed forward, the posterior covers $Q$, $Q_{\mathrm{kd}}$ and
the transcripts on a later horizon and under a knockdown that the data never contain. It misses
\pn{kim95.poolshift} of the 17 pools by about \pn{kim95.shiftlo}--\pn{kim95.shifthi} posterior
standard deviations, and that is where the posterior's error shows (\cref{sec:reducedpredict}).
Once every parameter is inferred, $Q$'s parameter variance is about \pn{kim96.ratio} times below the budget transferred from M0, the middle state of
\cref{sec:reducedpredict}. \todo{M10: how the posterior and these scores degrade as the data
become realistic (\cref{sec:ladder}).}

\paragraph{What carries and what does not.} $Q$ carries: a narrow interval, $\pm\pn{kim95.qint}\%$,
holds the truth at baseline and under the knockdown. The transcripts carry too, but about
two-thirds of their predictive variance is cell noise, so their hit fractions say little about
the parameter part of the prediction. The pools do not carry. For the trioses and FDP the
offsets come from $\theta$, not from cell noise, so they are biases of the posterior. The
knockdown's effect $\Delta Q$ is not resolved: the posterior shows only that it is small next
to one draw's cell noise. One reading of the pools is post hoc and untested. The promoter
strengths sit low in the posterior (\cref{sec:reduced}), yet the transcripts show no comparable
shift on the later horizon. Our reading is that the error the transcripts do not show moves
into metabolism, where the pools miss. We have not tested it.

\paragraph{The interface between modules.} The M0 analysis of \cref{sec:info} gives a working
rule. Single-cell data inform the parameters of their own module, and information between
modules travels on bulk observables and follows the gain of the coupling. A modeller can
compute the gains before any inference. On the reduced model the joint posterior uses both
streams at once and does not separate them, so the rule is not tested there. A dynamical reduction, such as the mean-field swap, changes the
simulator and keeps full Bayesian inference. A statistical cut keeps the simulator and changes
the target \cite{plummer2015cuts,jacob2017better,carmona2020semimodular}. The joint posterior
uses the first, and its cost was measured on M0. The modular arm's factor cut uses the second, and the joint posterior is the reference
against which its cost can be measured. \todo{KIM-116: whether
per-module inference lands close to the joint posterior, and at what fraction of its
CPU-hours.}

\paragraph{Limits.} Each of these limits names something we did not test.
\begin{itemize}
  \item \emph{One truth per design.} The reduced model has one dataset and one $\theta^\ast$.
  The truth inside \pn{kim92.truthin} of \pn{kim92.ncoord} intervals and the covered $Q$ are single
  events, not coverage rates. As in DREAM8 \cite{karr2015summary}, one known truth cannot show
  whether a posterior's width is right. Simulation-based calibration ran only on M0
  (\cref{sec:posteriors}).
  \item \emph{The exact sampler cannot check the reduced model.} It is projected at about
  \pn{cost.fullscale} CPU-hours per posterior, so on the reduced model the mixed route is tested
  only by its predictive check and by the scores against the true simulator.
  \item \emph{Synthetic data from the same model.} The model is well specified by construction.
  The \pn{kim96.sdpost}\% spread in $Q$ is precision under a correct model, and the pools show
  that it is not a calibrated uncertainty on every output. Real assays need observation models
  for unknown cell numbers in bulk metabolomics and for capture efficiency in single-cell RNA
  counts.
  \item \emph{The promoter shift.} Its mean standardised error is \pn{kim92.promz}, and the
  maximum a posteriori estimate shares it. One dataset cannot split it into noise that this
  dataset shares across promoters and a bias from fitting the mean-field bulk likelihood to data
  from the exact hybrid. A second dataset would.
  \item \emph{The gap on M0.} On one M0 dataset three parameters sit \pn{mix.gaplo}--\pn{mix.gaphi}
  standard deviations from the exact posterior, and we have not found the cause. \todo{KIM-90:
  the two leads, double counting of the transcripts' information through the bulk and the bulk
  factor itself, stated without a diagnosis.}
  \item \emph{A truth that is not sharp.} On $Q$, $Q_{\mathrm{kd}}$ and $\Delta Q$ the truth's
  standard error is \pn{kim95.sesdlo}--\pn{kim95.sesdhi} of the posterior standard deviation,
  above the sizing rule's \pn{kim95.sesdrule}.
  \item \emph{The one-way reduction.} The model keeps both directions of the coupling, but
  the mixed route treats the pools' effect on the rate constants as weak. In the scored runs GTP
  clips in \pn{kim95.clip}\% of baseline cells, at the truth and under the joint posterior alike,
  so it falls below $K_d$, where the small elasticities behind that treatment do not hold
  (\cref{sec:modelform}). The true simulator keeps the full coupling, so the scores include the
  effect of this reduction on the fit, and we have not measured it.
  \item \emph{The prior width.} The factor of \pn{kim96.ratio} holds at the default width only,
  and runs from \pn{kim96.widthlo} to \pn{kim96.widthhi} over widths $\tfrac12$ to 2. Many priors
  are asserted, with no width in their source.
  \item \emph{One organism, one reduction.} The framework is shown on the reduced \syn{} model
  only. Generality across organisms needs a second implementation.
\end{itemize}

\paragraph{What next.} The first step is real data, which needs the observation models above.
The second is the degradation ladder of \cref{sec:ladder}, which asks how far the generous
dataset can be degraded toward realistic assays before the posterior stops reaching its
parameters and the scores fall to the prior's. \todo{M10: what the ladder shows, and which assay
limits the posterior first.} The third is the full cell: the full well-stirred \syn{} model,
through the fast implementation of the companion paper \cite{kimpson2026whole}. There, as here, the
cost of a posterior is set by the number of simulations, not the time of one, so it depends on
the route more than on the simulator. A further question follows from the working rule: which parts of
a coupled simulator must be kept, and which latent variables inferred, to keep the information a
prediction needs? Each step adds reductions, and \cref{app:reductions} lists those of the reduced
model. Because each one is declared, a modeller who scores a prediction knows which departures
from the published model the score does not cover.
